% ============ 5. Alcance, riesgos y cronograma (D5) ============
\section{Alcance, riesgos y cronograma}

\subsection{Qué hará y qué no hará el proyecto}

\textbf{Sí hará:} auditar la calidad del padrón de investigadores reconocidos
(2013--2021) y de su artefacto consolidado; construir el panel longitudinal y
estimar transiciones y supervivencia; medir concentración territorial y
equidad de género y diversidad frente al censo DANE 2018; entregar un marco de
indicadores, un tablero de auditoría y documentación reproducible.

\textbf{No hará (y por qué):}
\begin{itemize}
  \item \textbf{No} validará el padrón contra registros internos de MinCiencias
        (no accesibles): se audita la fuente pública, que es el objeto
        declarado.
  \item \textbf{No} construirá modelos predictivos de carreras individuales ni
        análisis causal (requieren diseños fuera del alcance semestral y de los
        datos).
  \item \textbf{No} profundizará en el análisis de producción individual
        (índice $h$): exige el dataset de ~1,2\,GB y se declara como extensión.
  \item \textbf{No} incluirá convocatorias posteriores a 2021 (si 2023 se
        publicara a tiempo, se documenta como limitación, no se re-escopea).
\end{itemize}

\subsection{Riesgos principales y plan de contingencia}

\begin{enumerate}
  \item \textbf{Discrepancia de datos (77\,237 declarados vs.\ 50\,891 en el
        artefacto).} Ya detectada y caracterizada. Contingencia: se audita
        sobre el artefacto disponible y se documenta la discrepancia como
        hallazgo (es parte del objetivo, no un bloqueo).
  \item \textbf{Dataset de producción no descargable} (~1,2\,GB, disponibilidad
        intermitente). Contingencia: OE1--OE3 no dependen de él; se procesa si
        está disponible y se documenta la extensión.
  \item \textbf{Tiempo insuficiente para la inferencia completa.} Contingencia:
        se priorizan los entregables por hito y la holgura del cronograma (2
        semanas) absorbe retrasos puntuales.
\end{enumerate}

\subsection{Cronograma por hitos (12 semanas)}

\begin{longtable}{p{2.1cm}p{6.4cm}p{4.9cm}}
\caption{Cronograma por hitos con entregables verificables}
\label{tab:cronograma}\\
\toprule
\textbf{Semanas} & \textbf{Hito} & \textbf{Entregable verificable} \\
\midrule
\endfirsthead
\toprule
\textbf{Semanas} & \textbf{Hito} & \textbf{Entregable verificable} \\
\midrule
\endhead
1--2 & Datos abiertos y contratos & Scripts de ingesta/validación + catálogo actualizado (enlace y licencia verificados) \\
3--4 & Auditoría de calidad (OE1) & Tabla de indicadores de calidad por convocatoria + informe breve \\
5--6 & Panel longitudinal (OE2, parte 1) & Matrices de transición con IC y test de markovianidad \\
7--8 & Supervivencia (OE2, parte 2) & Curvas Kaplan--Meier + modelo Cox con covariables \\
9--10 & Equidad (OE3) & Indicadores HHI/Gini/Theil + razones de representación con IC \\
11 & Marco de indicadores y tablero & Tablero de auditoría + marco de indicadores \\
12 & Documentación y sustentación & Informe final reproducible + guion de sustentación \\
\midrule
\multicolumn{3}{l}{\emph{Holgura: semanas 6 y 11 incluyen margen de ~2 semanas
acumuladas para imprevistos.}} \\
\bottomrule
\end{longtable}

\textbf{Presupuesto de horas:} cada integrante dedicará \textbf{6 horas
semanales} (1 hora diaria después de clases, en la noche, 6 días a la semana)
$\times$ 3 integrantes $\times$ 12 semanas $\approx$ \textbf{216 horas}
disponibles; la carga estimada del plan es de \textbf{~180 horas}, dejando
\textbf{~36 horas (~17\,\%)} de holgura.
